\section*{Chapitre \ref{ch:g4:series-circuits} --- Les circuits en série}

\begin{solution}{exo:g4:series-circuits:1}
Les joueurs se trouvent sur une seule boucle, les uns après les
autres, sans jamais de séparation en deux routes --- comme des perles
enfilées sur un même collier.
\end{solution}

\begin{solution}{exo:g4:series-circuits:2}
Tout ou rien (une seule coupure arrête tout le monde)~; l'ordre est
sans importance (échanger les joueurs ne change rien)~; le partage
(plus il y a de lampes sur la boucle, plus chacune brille
faiblement).
\end{solution}

\begin{solution}{exo:g4:series-circuits:3}
La lampe 1 s'éteint elle aussi. La douille vide est une coupure dans
la boucle unique, et la règle première --- tout ou rien --- arrête
tous les joueurs d'un coup.
\end{solution}

\begin{solution}{exo:g4:series-circuits:4}
Rien ne change~: l'\omterm{ex:g3:first-electric-circuit:switch}{interrupteur} commande toujours les deux lampes
exactement comme avant. Règle deuxième --- l'ordre sur la boucle est
sans importance.
\end{solution}

\begin{solution}{exo:g4:series-circuits:5}
Deux lampes~: chacune nettement moins brillante que la lampe seule.
Trois~: plus faibles encore. C'est la poussée de la \omterm{def:g3:first-electric-circuit:battery}{pile} qui se
partage sur tout le collier.
\end{solution}

\begin{solution}{exo:g4:series-circuits:6}
Tête à queue~: le $+$ de l'une touchant le $-$ de la suivante, pour
que toutes deux poussent dans le même sens autour de la boucle.
Dessin~: $-\,[{+}\;{-}]\,[{+}\;{-}]\,+$ alignées sur une même file.
\end{solution}

\begin{solution}{exo:g4:series-circuits:7}
Depuis le $+$~: par la lame métallique du boîtier jusqu'au
\omterm{ex:g3:first-electric-circuit:switch}{bouton-interrupteur}, par l'\omterm{ex:g3:first-electric-circuit:switch}{interrupteur} jusqu'à la lampe, à travers
le filament de la lampe, puis retour le long du tube --- en
traversant la seconde \omterm{def:g3:first-electric-circuit:battery}{pile}, tête à queue --- jusqu'au $-$ de la
première \omterm{def:g3:first-electric-circuit:battery}{pile}. Une seule boucle, sans embranchement.
\end{solution}

\begin{solution}{exo:g4:series-circuits:8}
N'importe où sur la même boucle --- par exemple juste après le
premier bouton. \omterm{def:g4:series-circuits:series}{En série}, les deux boutons doivent être fermés pour
compléter l'unique route~: la sonnerie ne retentit donc que si l'on
appuie sur les deux en même temps.
\end{solution}

\begin{solution}{exo:g4:series-circuits:9}
Remplacer la première lampe par celle dont on sait qu'elle marche~;
si la guirlande s'allume, la lampe retirée était la coupable. Sinon,
remettre l'originale et passer à la deuxième douille, et ainsi de
suite tout au long de la guirlande. À une certaine douille, la
guirlande s'allumera --- la lampe qu'on vient de retirer est la perle
grillée.
\end{solution}

\begin{solution}{exo:g4:series-circuits:10}
Dans le noir, Tom a glissé une \omterm{def:g3:first-electric-circuit:battery}{pile} à l'envers. Tête contre tête, les
poussées des deux \omterm{def:g3:first-electric-circuit:battery}{piles} se combattent au lieu de s'ajouter, et la
lampe ne reçoit rien --- un match nul de pleine puissance reste un
match nul. Retourner une \omterm{def:g3:first-electric-circuit:battery}{pile} rétablit la file tête à queue, et la
\omterm{def:g1:light-and-shadows:source}{lumière}.
\end{solution}

\begin{solution}{exo:g4:series-circuits:11}
Elle a raison~: sur une boucle unique, un \omterm{ex:g3:first-electric-circuit:switch}{interrupteur} placé
n'importe où commande tout --- l'ordre est sans importance. Avec des
routes séparées, le courant pourrait contourner son \omterm{ex:g3:first-electric-circuit:switch}{interrupteur} par
une autre branche, et le moteur continuer de tourner~; les circuits à
branches demandent un placement d'\omterm{ex:g3:first-electric-circuit:switch}{interrupteur} plus réfléchi (ce sont
les circuits en dérivation de l'an prochain).
\end{solution}
